\chapter{Observability: Tracing, Cost Accounting and Logging}
\label{ch:observability}

Every answer VidyaRAG returns carries a record of how it was produced: how long each stage took, how many tokens were used, what that would cost, and what the guards, reranker and self-check did. This chapter explains that record, where it surfaces, and two measurement defects found in it during this analysis: time counted twice, and tokens not counted at all.

\section{Design intent}

\enquote{The trace is returned with the answer rather than only logged, so latency and cost are inspectable from the API and the UI, not just from a log file nobody opens} \ev{src/vidyarag/observe/trace.py:3-7}. On cost, the same docstring explains why a list price is shown although the free tier costs nothing: reporting \enquote{\$0.00/query} would \enquote{hide the thing that actually matters --- whether the design is economical enough to run for real}.

\section{The trace object}

\begin{lstlisting}[language=Python]
class QueryTrace(BaseModel):
    query: str
    profile: str = "baseline"
    prompt_version: str = ""
    stages: list[StageTiming] = []          # (name, duration_ms), in the order recorded
    usage: list[Usage] = []                 # (model, input_tokens, output_tokens, purpose)
    retrieved_chunk_ids: list[str] = []     # candidate pool before reranking
    ranked_chunk_ids: list[str] = []        # final ordering
    cited_chunk_ids: list[str] = []
    attempts: int = 1
    abstained: bool = False
    sub_questions: list[str] = []
    corrective: dict[str, object] = {}      # loop outcome, one entry per attempt
    guard_input: dict[str, object] = {}
    guard_context: dict[str, object] = {}
    rerank: dict[str, object] = {}          # reordered / promoted_into_context / top1_changed

    @property
    def total_ms(self) -> float:
        return sum(stage.duration_ms for stage in self.stages)

    @contextmanager
    def stage(self, name: str) -> Iterator[None]:
        started = time.perf_counter()
        try:
            yield
        finally:
            self.record(name, (time.perf_counter() - started) * 1000)
\end{lstlisting}
\vspace{-4pt}{\raggedleft\ev{src/vidyarag/observe/trace.py:75-179} (docstrings removed, defaults simplified)\par}

A stage is recorded in \code{finally}, so \enquote{a failed query still shows where the time went --- which is exactly when that matters} \ev{src/vidyarag/observe/trace.py:168-179}. The module also defines a \code{Timer} class, which nothing in the repository uses.

\subsection{Which code records which stage}

\begin{center}
\begin{tabularx}{\textwidth}{@{}>{\raggedright\arraybackslash}p{2.6cm}>{\raggedright\arraybackslash}p{5.2cm}L@{}}
\toprule
\textbf{Stage} & \textbf{Recorded in} & \textbf{Covers}\\
\midrule
\code{guard_input} & \code{pipeline.py:160} & Input screening\\
\code{decompose} & \code{pipeline.py:93} & The Gemini decomposition call\\
\code{retrieve} & \code{pipeline.py:98} & Query embedding and vector search (one per sub-question when decomposed)\\
\code{rerank} & \code{pipeline.py:121} & Cross-encoder scoring of the candidates\\
\code{guard_context} & \code{pipeline.py:137} & Passage screening\\
\code{generate} & \code{generate/answer.py:106} & The Gemini answer call only\\
\code{corrective} & \code{pipeline.py:230} & The whole self-check loop, \emph{including} every stage above that runs inside it, plus grading\\
\bottomrule
\end{tabularx}
\end{center}

\section{Token usage and list price}

\begin{lstlisting}[language=Python]
GEMINI_PRICING: dict[str, tuple[float, float]] = {   # USD per million tokens, "verified August 2026"
    "gemini-3.5-flash": (0.30, 2.50),
    "gemini-3.5-flash-lite": (0.10, 0.40),
    "gemini-3.1-flash-lite": (0.10, 0.40),
    "gemini-3.6-flash": (0.30, 2.50),
}
_FALLBACK_PRICING = (0.30, 2.50)

def price_call(model: str, input_tokens: int, output_tokens: int) -> float:
    rate_in, rate_out = GEMINI_PRICING.get(model, _FALLBACK_PRICING)
    return (input_tokens * rate_in + output_tokens * rate_out) / 1_000_000
\end{lstlisting}
\vspace{-4pt}{\raggedleft\ev{src/vidyarag/observe/trace.py:26-44}\par}

For a call with $t_{\text{in}}$ input and $t_{\text{out}}$ output tokens at rates $r_{\text{in}}, r_{\text{out}}$ (USD per million):
\begin{equation}
\text{price} = \frac{t_{\text{in}}\, r_{\text{in}} + t_{\text{out}}\, r_{\text{out}}}{10^6}.
\end{equation}
An unknown model is priced at Flash rates rather than zero: \enquote{a silently free-looking call is worse than a slightly wrong estimate.} The rates are written in the code; whether they match Google's current price list is \NotDet

\begin{example}[the shipped profile's reported cost per question]
The committed \code{guarded} run averages 2,423 input and 66 output tokens per question on \code{gemini-3.5-flash-lite}: $(2{,}423 \times 0.10 + 66 \times 0.40)/10^6 = (242.3 + 26.4)/10^6 \approx \$0.00027$, the per-query figure in the README. Section~\ref{sec:token-bug} explains why this figure is too low.
\end{example}

\section{Where the trace surfaces}

\begin{center}
\begin{tabularx}{\textwidth}{@{}>{\raggedright\arraybackslash}p{3.6cm}LLLL@{}}
\toprule
\textbf{Trace content} & \textbf{API \code{TraceOut}} & \textbf{Demo panel} & \textbf{CLI \code{ask}} & \textbf{Evaluation run file}\\
\midrule
Stage timings & \yes & \yes & \yes (summary line) & total only\\
Total ms & \yes & \yes & \yes & \yes (\code{latency_ms})\\
Tokens, list price & \yes & \yes & \yes & \yes\\
Retrieved / cited counts & \yes & retrieved only & --- & ids\\
Abstained flag & \no & as an event & --- & \yes\\
Corrective attempts and scores & \no & attempt count & --- & \yes\\
Guard events & \no & \yes & --- & ---\\
Rerank movement, sub-questions & \no & \no & --- & \yes\\
\bottomrule
\end{tabularx}
\end{center}
\vspace{-4pt}{\raggedleft\ev{src/vidyarag/api/models.py:50-66}, \ev{app/app.py:104-144}, \ev{src/vidyarag/cli.py:271}\par}

The README quickstart shows a CLI summary line of this form: \code{10437ms [retrieve=3163ms generate=7274ms] 2553+227 tok}. The \code{summary} method appends the list price in parentheses \ev{src/vidyarag/observe/trace.py:181-188}.

\section{Defect 1: nested stages are counted twice}
\label{sec:latency-bug}

\subsection{Mechanism}

\begin{lstlisting}[language=Python]
# pipeline.py, _answer_corrective
with trace.stage("corrective"):                  # outer stage...
    outcome = run_corrective_loop(
        question=question,
        generate=generate,    # -> generate_answer -> trace.stage("generate")
        retrieve=retrieve,    # -> self.retrieve   -> trace.stage("retrieve"), ("rerank"), ("guard_context")
        ...)

# trace.py
@property
def total_ms(self) -> float:
    return sum(stage.duration_ms for stage in self.stages)   # ...and every inner stage, again
\end{lstlisting}
\vspace{-4pt}{\raggedleft\ev{src/vidyarag/pipeline.py:213-238}, \ev{src/vidyarag/observe/trace.py:128-130}\par}

The \code{corrective} stage's duration already contains the retrieval, reranking, screening and generation stages that run inside it. \code{total_ms} adds them all, so for the \code{corrective} and \code{guarded} profiles:
\begin{equation}
\text{reported } \texttt{total\_ms} \;=\; \underbrace{t_{\text{guard\_input}} + t_{\text{corrective}}}_{\text{real elapsed time}} \;+\; \underbrace{t_{\text{retrieve}} + t_{\text{rerank}} + t_{\text{guard\_context}} + t_{\text{generate}}}_{\text{counted a second time}}.
\end{equation}
The other three profiles have no enclosing stage and are unaffected. Retries sum correctly: a second retrieval really does take more time, and a test asserts that a stage timed twice counts twice (\code{test_same_stage_can_be_timed_more_than_once}). No test checks nesting, or compares \code{total_ms} with wall-clock time.

\subsection{Evidence}

\begin{measured}[three \code{guarded} queries, timed around \code{Pipeline.answer} on the development machine]
\begin{center}\small
\begin{tabular}{@{}lrrrr@{}}
\toprule
\textbf{Question} & \textbf{Wall (ms)} & \textbf{Reported (ms)} & \textbf{Nested (ms)} & \textbf{Excess (ms)}\\
\midrule
Glucose, facilitated diffusion & 79,673 & 92,578 & 12,906 & 12,905\\
Oxytocin thresholds (abstained) & 16,463 & 29,060 & 12,596 & 12,597\\
Anaphase of mitosis & 30,848 & 43,474 & 12,626 & 12,626\\
\bottomrule
\end{tabular}
\end{center}
In every case the excess over wall-clock time equals the sum of the nested stages to within a millisecond. (The first query spent about 67 seconds in the grading call alone, a free-tier latency spike.) On the live Space, one \code{guarded} answer reported a total of 9,746\,ms while the round trip measured from the client, network included, was 7.1 seconds --- impossible unless time is counted twice.
\end{measured}

\begin{figure}[htbp]
\centering
\resizebox{\textwidth}{!}{%
\begin{tikzpicture}[font=\sffamily\scriptsize, x=0.17cm, y=0.6cm]
\draw[-Stealth] (0,0) -- (97,0) node[right] {seconds};
\foreach \t in {0,10,20,30,40,50,60,70,80,90} { \draw (\t,0.12) -- (\t,-0.12) node[below] {\t}; }
\fill[blue!25] (0,5) rectangle (79.67,5.7);
\node[anchor=west] at (0,6.1) {\code{corrective} stage (recorded): 79,668 ms};
\fill[orange!45] (0,3.9) rectangle (12.91,4.6);
\node[anchor=west] at (13.4,4.25) {inside it, also recorded: retrieve 449 + rerank 10,538 + guard\_context 17 + generate 1,902 = 12,906 ms};
\fill[gray!25] (12.91,2.8) rectangle (79.67,3.5);
\node at (46.3,3.15) {grading call --- inside \code{corrective}, not a stage of its own};
\draw[very thick, green!45!black] (0,1.9) -- (79.67,1.9);
\node[anchor=west] at (80.3,1.9) {wall clock 79,673 ms};
\draw[very thick, red!70!black] (0,1.0) -- (92.58,1.0);
\node[anchor=south east] at (92.58,1.05) {reported \code{total_ms} 92,578 ms = 79,668 + 12,906 (+ guard\_input)};
\end{tikzpicture}}
\caption{The first measured query on a time axis. The orange stages are inside the blue one, but \code{total_ms} adds both.}
\label{fig:latency-bug}
\end{figure}

\subsection{Consequences}

\begin{itemize}
  \item \textbf{Committed latency for the self-check profiles is overstated.} The evaluation runner stores \code{latency_ms = trace.total_ms} for each question \ev{src/vidyarag/evaluation/runner.py:291}, so the reported means of 22,031\,ms (\code{corrective}) and 21,738\,ms (\code{guarded}) include the doubled stages. The \code{baseline} (1,091\,ms), \code{rerank} (6,856\,ms) and \code{decompose} (2,635\,ms) figures are not affected. \tagV
  \item \textbf{The demo's \enquote{total} row is overstated} for every self-check answer, because it prints the same property \ev{app/app.py:111}. So is the API's \code{trace.total_ms}. \tagV
  \item \textbf{Statements built on the number inherit the error}, including \code{docs/DESIGN.md}'s \enquote{22 s per query against the baseline's 1.1 s}. \tagV
\end{itemize}

\begin{inferred}[how large the error is in the committed runs]
The run files store only each question's total, not its stages, so the true mean cannot be recomputed from them. The doubled part is roughly the retrieval, reranking and generation time --- about what the \code{rerank} profile takes on its own, whose mean is 6,856\,ms and median 3,100\,ms. The true \code{guarded} mean is therefore probably between about 15 and 19 seconds rather than 21.7. Re-running the evaluation after a fix is the only way to know.
\end{inferred}

\section{Defect 2: grading and decomposition tokens are never recorded}
\label{sec:token-bug}

\subsection{Mechanism}

\code{QueryTrace.add_usage} is called from exactly one place, \code{generate/answer.py} \ev{src/vidyarag/generate/answer.py:62-70}. \code{grade_answer} and \code{decompose} both call Gemini and both receive usage metadata, and neither records it. The \code{Usage.purpose} field's docstring anticipates \enquote{generation, grading, decomposition} \ev{src/vidyarag/observe/trace.py:63-64}; only \code{generation} is ever written.

\subsection{Evidence}

\begin{measured}[the unrecorded grading prompt, for the same three queries]
The recorded generation inputs were 2,554, 2,392 and 2,367 tokens. The grading prompts for the same answers, measured with Gemini's token-counting endpoint, were 2,578, 2,336 and 2,471 tokens, none of which appears in the trace. In a separate measurement on the anaphase answer (Chapter~\ref{ch:trace}), a successful grading call used 2,497 input and 612 output tokens, against 2,367 and 156 for the generation call. At the rates in the code, the unrecorded grading call cost more than the recorded generation call: about \$0.00049 against \$0.00030.
\end{measured}

\subsection{Consequences}

\begin{itemize}
  \item For \code{corrective} and \code{guarded}, every graded question makes a second model call whose input is about the same size as generation's. Reported input tokens, and therefore the list price, are understated by roughly half; the grading output is missing on top of that. \tagI
  \item The README's \$0.00027 per query for the shipped profile, and the demo's \enquote{Cost at list price}, are too low for the same reason. \tagV
  \item The \code{decompose} profile's figures omit the decomposition call. \tagV
\end{itemize}

\begin{interview}
These two defects are good material, not embarrassments, if you present them as findings: \enquote{My trace summed nested spans, so the self-check profiles over-reported latency by the retrieval and generation time; I proved it by timing three queries --- the excess matched the nested stages to the millisecond. Separately, only the generator recorded tokens, so cost was understated by about half for the self-check.} Then say how you would fix and guard each: a wall-clock total, or parent-aware spans; usage recorded at every model call; and a test for each.
\end{interview}

\section{Logging}

\begin{lstlisting}[language=Python]
def configure_logging(level: str = "INFO", fmt: str = "json") -> None:
    logging.basicConfig(format="%(message)s", stream=sys.stdout,
                        level=getattr(logging, level.upper(), logging.INFO))
    renderer = structlog.dev.ConsoleRenderer() if fmt == "console" \
               else structlog.processors.JSONRenderer()
    structlog.configure(
        processors=[structlog.contextvars.merge_contextvars,
                    structlog.processors.add_log_level,
                    structlog.processors.TimeStamper(fmt="iso", utc=True),
                    structlog.processors.StackInfoRenderer(),
                    structlog.processors.format_exc_info,
                    renderer],
        wrapper_class=structlog.make_filtering_bound_logger(
            getattr(logging, level.upper(), logging.INFO)),
        cache_logger_on_first_use=True)
\end{lstlisting}
\vspace{-4pt}{\raggedleft\ev{src/vidyarag/observe/logging.py:19-52} (abridged)\par}

The only callers are in the API's startup code, which logs two events: \code{api.startup} (with profile, collection, embedding model and generation model) and \code{api.shutdown} \ev{src/vidyarag/api/main.py:30-50}. Nothing logs a query, a guard event, an abstention, a grader failure or an error. The CLI and the demo never configure structlog.

\section{What is not implemented}

\begin{itemize}
  \item No metrics endpoint or time series (for example request counts, latency percentiles, abstention rate).
  \item No distributed tracing or request identifiers.
  \item No persisted query log or analytics.
  \item No error tracking or alerting --- in particular, nothing reports a run of unmeasured (failed) grading calls, during which the self-check accepts answers unverified (Chapter~\ref{ch:corrective}).
\end{itemize}

\begin{recommended}
Compute the total as wall-clock time around \code{Pipeline.answer}, or record the loop's inner stages as children and exclude them from the sum; add a test that a self-check query's total is within a few milliseconds of its wall time. Record usage in \code{grade_answer} and \code{decompose} with purposes \code{grading} and \code{decomposition}, then re-run the evaluation to correct the latency and cost figures. Emit one structured log line per query with the trace summary, the abstention flag and any guard or grader error. Add \code{abstained}, the attempt count and guard events to the API's \code{TraceOut}.
\end{recommended}
